\documentclass{article}
\usepackage[utf8]{inputenc}
\usepackage[margin=0.5in]{geometry}
\usepackage{booktabs}
\usepackage{longtable}

\title{PCD Homework 1 Report}
\author{Vasilică Alexandru, MISS1}
\begin{document}
\maketitle

\section{Experiment Description}
This experiment evaluates the performance and reliability of five different transport protocol implementations:
\begin{itemize}
    \item \textbf{TCP}: Standard stream-oriented reliable transport.
    \item \textbf{UDP}: Raw unreliable datagram transport.
    \item \textbf{UDP (Stop and Wait)}: A custom reliability layer implemented over UDP requiring an ACK for every packet.
    \item \textbf{QUIC}: Used in Datagram mode to allow for high-throughput and to allow for testing the protocol's data loss without the built in reliability of stream mode.
    \item \textbf{QUIC (Stop and Wait)}: A custom reliability layer implemented over QUIC datagrams.
\end{itemize}

The metrics were collected by transferring two different payload sizes (500 MB and 1 GB) using varying block sizes (512 B, 1KB, 16KB, ~60KB) over a local loopback connection. Quic does not support datagram sizes larger than 1.5KB, so the experiments were run with 512 B and 1KB.

\section{How to Run}
The project requires Rust and Cargo to be installed. The server and client can be run using the following commands:
\subsection{Running the Server}

The server listens to traffic for all the protocols concurrently on separate ports. To start the server, run the following command:

\begin{verbatim}
cargo run -p server
\end{verbatim}

\subsection{Running the Client}
\begin{verbatim}
cargo run -p client -- --transport <TYPE> --file-path <PATH> --block-size <BYTES>
\end{verbatim}

\subsubsection{Client Options}
\begin{itemize}
    \item \textbf{--transport TYPE}: The transport protocol to use. Can be \textbf{tcp}, \textbf{udp}, \textbf{udp-stop-and-wait}, \textbf{quic}, or \textbf{quic-stop-and-wait}.
    \item \textbf{--file-path PATH}: The path to the file to transfer.
    \item \textbf{--block-size BYTES}: The size of the block to transfer.
\end{itemize}

Note: When using UDP with larger block sizes, the UDP maximum datagram size of the system may need to be increased.
\begin{verbatim}
sudo sysctl -w net.inet.udp.maxdgram=65535
\end{verbatim}

\section{Experiment Results}

\subsection{ TCP}
\begin{small}
    \begin{longtable}{l r r r r r r r}
        \toprule
        \textbf{Payload} & \textbf{Block} & \textbf{Time(s)} & \textbf{Pkts(S)} & \textbf{Pkts(R)} & \textbf{Bytes(S)} & \textbf{Bytes(R)} & \textbf{Loss\%} \\
        \midrule
        1 GB             & 512            & 19.6306          & 2097152          & 2097152          & 1073741824        & 1073741824        & 0.00\%          \\
        1 GB             & 1024           & 9.6824           & 1048576          & 1048576          & 1073741824        & 1073741824        & 0.00\%          \\
        1 GB             & 16384          & 1.2510           & 65536            & 65536            & 1073741824        & 1073741824        & 0.00\%          \\
        1 GB             & 60000          & 0.5400           & 17896            & 17896            & 1073741824        & 1073741824        & 0.00\%          \\
        500 MB           & 512            & 10.4824          & 1024000          & 1024000          & 524288000         & 524288000         & 0.00\%          \\
        500 MB           & 1024           & 5.0572           & 512000           & 512000           & 524288000         & 524288000         & 0.00\%          \\
        500 MB           & 16384          & 0.6705           & 32000            & 32000            & 524288000         & 524288000         & 0.00\%          \\
        500 MB           & 60000          & 0.2905           & 8739             & 8739             & 524288000         & 524288000         & 0.00\%          \\
        \bottomrule
    \end{longtable}
\end{small}
\subsection{ UDP}
\begin{small}
    \begin{longtable}{l r r r r r r r}
        \toprule
        \textbf{Payload} & \textbf{Block} & \textbf{Time(s)} & \textbf{Pkts(S)} & \textbf{Pkts(R)} & \textbf{Bytes(S)} & \textbf{Bytes(R)} & \textbf{Loss\%} \\
        \midrule
        1 GB             & 512            & 20.0796          & 2097152          & 2094517          & 1073741824        & 1072392704        & 0.13\%          \\
        1 GB             & 1024           & 10.1761          & 1048576          & 1046042          & 1073741824        & 1071147008        & 0.24\%          \\
        1 GB             & 16384          & 0.7629           & 65536            & 62467            & 1073741824        & 1023459328        & 4.68\%          \\
        1 GB             & 60000          & 0.3009           & 17896            & 15836            & 1073741824        & 950141824         & 11.51\%         \\
        500 MB           & 512            & 9.8073           & 1024000          & 1021243          & 524288000         & 522876416         & 0.27\%          \\
        500 MB           & 1024           & 4.9123           & 512000           & 509365           & 524288000         & 521589760         & 0.51\%          \\
        500 MB           & 16384          & 0.3790           & 32000            & 28771            & 524288000         & 471384064         & 10.09\%         \\
        500 MB           & 60000          & 0.1666           & 8739             & 7062             & 524288000         & 423668000         & 19.19\%         \\
        \bottomrule
    \end{longtable}
\end{small}
\subsection{ UDP (Stop and Wait)}
\begin{small}
    \begin{longtable}{l r r r r r r r}
        \toprule
        \textbf{Payload} & \textbf{Block} & \textbf{Time(s)} & \textbf{Pkts(S)} & \textbf{Pkts(R)} & \textbf{Bytes(S)} & \textbf{Bytes(R)} & \textbf{Loss\%} \\
        \midrule
        1 GB             & 512            & 110.0781         & 2097152          & 2097152          & 1073741824        & 1073741824        & 0.00\%          \\
        1 GB             & 1024           & 55.2263          & 1048576          & 1048576          & 1073741824        & 1073741824        & 0.00\%          \\
        1 GB             & 16384          & 3.6949           & 65536            & 65536            & 1073741824        & 1073741824        & 0.00\%          \\
        1 GB             & 60000          & 1.0859           & 17896            & 17896            & 1073741824        & 1073741824        & 0.00\%          \\
        500 MB           & 512            & 53.7394          & 1024000          & 1024000          & 524288000         & 524288000         & 0.00\%          \\
        500 MB           & 1024           & 26.9609          & 512000           & 512000           & 524288000         & 524288000         & 0.00\%          \\
        500 MB           & 16384          & 1.7910           & 32000            & 32000            & 524288000         & 524288000         & 0.00\%          \\
        500 MB           & 60000          & 0.5494           & 8739             & 8739             & 524288000         & 524288000         & 0.00\%          \\
        \bottomrule
    \end{longtable}
\end{small}

\subsection{ QUIC}
\begin{small}
    \begin{longtable}{l r r r r r r r}
        \toprule
        \textbf{Payload} & \textbf{Block} & \textbf{Time(s)} & \textbf{Pkts(S)} & \textbf{Pkts(R)} & \textbf{Bytes(S)} & \textbf{Bytes(R)} & \textbf{Loss\%} \\
        \midrule
        1 GB             & 512            & 14.6380          & 2097152          & 2041553          & 1073741824        & 1045275136        & 2.65\%          \\
        1 GB             & 1024           & 13.4797          & 1048576          & 989332           & 1073741824        & 1013075968        & 5.65\%          \\
        500 MB           & 512            & 6.7652           & 1024000          & 973617           & 524288000         & 498491904         & 4.92\%          \\
        500 MB           & 1024           & 5.9717           & 512000           & 447697           & 524288000         & 458441728         & 12.56\%         \\
        \bottomrule
    \end{longtable}
\end{small}
\subsection{ QUIC (Stop and Wait)}
\begin{small}
    \begin{longtable}{l r r r r r r r}
        \toprule
        \textbf{Payload} & \textbf{Block} & \textbf{Time(s)} & \textbf{Pkts(S)} & \textbf{Pkts(R)} & \textbf{Bytes(S)} & \textbf{Bytes(R)} & \textbf{Loss\%} \\
        \midrule
        1 GB             & 512            & 151.8875         & 2097152          & 2097152          & 1073741824        & 1073741824        & 0.00\%          \\
        1 GB             & 1024           & 75.7086          & 1048576          & 1048576          & 1073741824        & 1073741824        & 0.00\%          \\
        500 MB           & 512            & 74.0261          & 1024000          & 1024000          & 524288000         & 524288000         & 0.00\%          \\
        500 MB           & 1024           & 37.0389          & 512000           & 512000           & 524288000         & 524288000         & 0.00\%          \\
        \bottomrule
    \end{longtable}
\end{small}

\section{Results Analysis}
\begin{itemize}
    \item \textbf{Throughput Efficiency}: For all protocols, increasing the block size resulted in a massive reduction in transfer time due to lower per-packet overhead and fewer system calls.
    \item \textbf{Reliability vs. Latency}: The "Stop and Wait" implementations achieved 100\% reliability but were significantly slower as they wait for an acknowledgement for every individual packet.
    \item \textbf{Block Size vs. Data Loss}: In unreliable transports (UDP and QUIC), there is a direct correlation between block size and the probability of data loss. Larger blocks (e.g., 60,000 B) force the underlying layers to use IP fragmentation (in UDP). Because datagrams are atomic, the loss of a single byte or fragment results in the discarding of the entire block. Consequently, while larger blocks improve throughput by reducing header-to-payload ratios, they significantly increase the loss risk, leading to the observed 11--19\% loss rates compared to the sub-1\% loss seen with smaller blocks.

    \item \textbf{Security Overhead}: The mandatory TLS 1.3 encryption in QUIC introduces a measurable performance "tax" compared to raw UDP. Each datagram requires CPU-intensive cryptographic operations for packet protection.

    \item \textbf{Time Efficiency}: Raw UDP was the fastest overall for large blocks (e.g., 0.30s for 1 GB), as it has the lowest per-packet overhead and no congestion control. TCP followed closely (0.54s for 1 GB) due to its highly optimized OS-level stack. Interestingly, at small block sizes (512 B), QUIC (14.64s) was more efficient than both TCP (19.63s) and UDP (20.08s), demonstrating QUIC's superior handling of high-packet-rate streams and its reduced handshake overhead for the first few packets.
\end{itemize}

\section{Conclusions}
The experiments demonstrate that there is no "one-size-fits-all" protocol; the choice depends heavily on the specific requirements for reliability, latency, and security.

\subsection{TCP (Transmission Control Protocol)}
\begin{itemize}
    \item \textbf{Advantages}: Guaranteed 100\% reliability and ordering. Highly optimized OS-level implementation. No data loss regardless of block size.
    \item \textbf{Disadvantages}: Vulnerable to "Head-of-Line Blocking". Higher overhead due to heavy handshake and congestion control.

\end{itemize}

\subsection{UDP (User Datagram Protocol)}
\begin{itemize}
    \item \textbf{Advantages}: Minimal overhead and lowest possible latency. Maximum theoretical throughput for large block sizes.
    \item \textbf{Disadvantages}: No reliability or ordering guarantees. High packet loss for large blocks due to IP fragmentation sensitivity.
\end{itemize}

\subsection{QUIC (Quick UDP Internet Connections)}
\begin{itemize}
    \item \textbf{Advantages}: Built-in TLS 1.3 security. Better handling of high-packet-rate streams (small blocks). Supports connection migration.
    \item \textbf{Disadvantages}: Measurable CPU overhead from mandatory encryption. Datagram mode is strictly MTU-limited and lacks native fragmentation.

\end{itemize}



\end{document}